\documentclass[11pt,a4paper]{article}
\usepackage[T1]{fontenc}
\usepackage[utf8]{inputenc}
\usepackage{lmodern}
\usepackage{geometry}
\usepackage{amsmath,amssymb}
\usepackage{amsthm}
\usepackage{mathpartir}
\usepackage{listings}
\usepackage{xcolor}
\geometry{margin=25mm}

\newtheorem{theorem}{Theorem}
\newtheorem{lemma}[theorem]{Lemma}
\newtheorem{remark}[theorem]{Remark}

\lstdefinelanguage{Coq}{
  morekeywords={
    Inductive,Definition,Fixpoint,Lemma,Theorem,Proof,Qed,forall,exists,
    match,with,end,fun,Prop,Type,Set,apply,eapply,constructor,intros,
    inversion,subst,destruct,reflexivity,exact
  },
  sensitive=true,
  morecomment=[n]{(*}{*)}
}
\lstset{
  language=Coq,
  basicstyle=\ttfamily\small,
  keywordstyle=\color{blue!60!black},
  commentstyle=\color{green!35!black},
  breaklines=true,
  frame=single,
  columns=fullflexible
}

\newcommand{\tint}{\mathtt{int}}
\newcommand{\inst}{\preceq}
\newcommand{\gen}{\mathrm{gen}}
\newcommand{\ftv}{\mathrm{ftv}}
\newcommand{\subst}[3]{#1[#2 := #3]}
\newcommand{\step}{\longrightarrow}

\title{Coq Proofs for Hindley--Milner Inference Soundness\\and Core Type Safety}
\author{}
\date{}

\begin{document}
\maketitle

\begin{abstract}
This is an abridged English rendering of \texttt{hm\_safety\_report\_ja.tex},
which is the primary version and carries the proofs step by step; the proof
sketches below are shorter.  It describes the two Coq developments in this
repository:
\texttt{HMSoundness.v}, which relates an algorithmic Hindley--Milner inference
judgement to a declarative one, and \texttt{HMTypeSafety.v}, which proves
progress and preservation for a typed core language.  Both are set out here in
the usual mathematical notation, with the grammars, the type definitions and
the inference rules written out, so that the Coq text can be read against what
it encodes, and with the shape of each proof described rather than merely
cited.  Section~\ref{sec:reading} states plainly what the two theorems do and
do not establish; that is the section to read before quoting any result from
here.
\end{abstract}

\tableofcontents

\section{Goal}

\begin{enumerate}
  \item \emph{Inference soundness}: if the inference judgement assigns a type
  to an expression, the declarative judgement assigns it too.  This is
  \texttt{HMSoundness.v}.
  \item \emph{Type safety}: a closed well-typed term is a value or can take a
  step, and stepping preserves its type.  This is \texttt{HMTypeSafety.v}.
\end{enumerate}

The two files are independent.  Each defines its own syntax, its own types and
its own judgements, and neither imports the other.  In particular the type
language of the second is not the type language of the first, a point taken up
in Section~\ref{sec:reading}.

% =====================================================================
\section{Inference soundness: syntax}
\label{sec:sound-syntax}

\subsection{Grammar}

\[
\begin{array}{lcll}
  e & ::= & x \mid n \mid \lambda x.\,e \mid e_1\,e_2
            \mid \mathtt{let}\;x = e_1\;\mathtt{in}\;e_2
    & \text{expressions}\\[2pt]
  \tau & ::= & \tint \mid \alpha \mid \tau_1 \to \tau_2
    & \text{types}\\[2pt]
  \sigma & ::= & \tau \mid \forall \bar\alpha.\,\tau
    & \text{type schemes}
\end{array}
\]

Variables $x$ and type variables $\alpha$ are natural numbers, and
$\bar\alpha$ is a list of them.  In Coq:

\begin{lstlisting}
Inductive ty : Type :=
| TInt : ty
| TVar : nat -> ty
| TArrow : ty -> ty -> ty.

Inductive expr : Type :=
| EVar : nat -> expr   | EInt : nat -> expr
| ELam : nat -> expr -> expr
| EApp : expr -> expr -> expr
| ELet : nat -> expr -> expr -> expr.

Inductive scheme : Type :=
| Mono : ty -> scheme
| Forall : list nat -> ty -> scheme.
\end{lstlisting}

A scheme is either monomorphic or a $\forall$ wrapper; these are two separate
constructors rather than the wrapper with an empty list.

\subsection{Environments}

An environment is a partial map from variables to schemes, represented as a
function.
\[
  \Gamma : \mathbb{N} \to \sigma_\bot,
  \qquad
  \emptyset = \lambda y.\,\bot,
  \qquad
  (\Gamma, x{:}\sigma)(y) =
    \begin{cases}
      \sigma & y = x\\
      \Gamma(y) & y \neq x
    \end{cases}
\]

\begin{lstlisting}
Definition env := nat -> option scheme.
Definition empty : env := fun _ => None.
Definition extend (G : env) (x : nat) (s : scheme) : env :=
  fun y => if Nat.eqb x y then Some s else G y.
\end{lstlisting}

Because environments are functions rather than lists, $\ftv(\Gamma)$ is not
computable here.  That is why the generalization rule below carries no side
condition.

\subsection{Instantiation and generalization}

\begin{mathpar}
  \inferrule*[right=Inst-Mono]
    { }
    {\tau \;\inst\; \tau}
  \and
  \inferrule*[right=Inst-Forall]
    { }
    {\forall\bar\alpha.\,\tau \;\inst\; \tau}
  \and
  \inferrule*[right=Gen-Mono]
    { }
    {\gen(\Gamma, \tau) \ni \tau}
  \and
  \inferrule*[right=Gen-Forall]
    { }
    {\gen(\Gamma, \tau) \ni \forall\bar\alpha.\,\tau}
\end{mathpar}

These are the rules as the file states them, and they repay careful reading.

\begin{remark}[Instantiation does not substitute]
\label{rem:inst}
\textsc{Inst-Forall} relates $\forall\bar\alpha.\,\tau$ to $\tau$ and to
nothing else.  A scheme therefore has exactly one instance, its own body: the
quantified variables are never replaced.  Instantiation is the identity on
types.
\end{remark}

\begin{remark}[Generalization has no side condition]
\label{rem:gen}
\textsc{Gen-Forall} holds for \emph{every} list $\bar\alpha$, with no
requirement that those variables be absent from $\Gamma$.  The Damas--Milner
condition $\bar\alpha \cap \ftv(\Gamma) = \emptyset$ is not present.  Nothing
goes wrong as a consequence, but only because Remark~\ref{rem:inst} means the
quantifier is never eliminated: a scheme built by \textsc{Gen-Forall} can only
ever be used at its own body.
\end{remark}

The file's own header says as much: instantiation and generalization are
``kept as relations \dots\ not on a particular implementation of free-variable
sets and substitution''.  The consequences for the main theorem are drawn in
Section~\ref{sec:reading}.

\section{Inference soundness: rules}

\subsection{Declarative system}

\begin{mathpar}
  \inferrule*[right=T-Var]
    {\Gamma(x) = \sigma \\ \sigma \inst \tau}
    {\Gamma \vdash x : \tau}
  \and
  \inferrule*[right=T-Int]
    { }
    {\Gamma \vdash n : \tint}
  \and
  \inferrule*[right=T-Lam]
    {\Gamma, x{:}\tau_1 \vdash e : \tau_2}
    {\Gamma \vdash \lambda x.\,e : \tau_1 \to \tau_2}
  \and
  \inferrule*[right=T-App]
    {\Gamma \vdash e_1 : \tau_1 \to \tau_2 \\ \Gamma \vdash e_2 : \tau_1}
    {\Gamma \vdash e_1\,e_2 : \tau_2}
  \and
  \inferrule*[right=T-Let]
    {\Gamma \vdash e_1 : \tau_1 \\
     \gen(\Gamma,\tau_1) \ni \sigma \\
     \Gamma, x{:}\sigma \vdash e_2 : \tau_2}
    {\Gamma \vdash \mathtt{let}\;x = e_1\;\mathtt{in}\;e_2 : \tau_2}
\end{mathpar}

\subsection{Algorithmic system}

\begin{mathpar}
  \inferrule*[right=I-Var]
    {\Gamma(x) = \sigma \\ \sigma \inst \tau}
    {\Gamma \vdash_{\!I} x : \tau}
  \and
  \inferrule*[right=I-Int]
    { }
    {\Gamma \vdash_{\!I} n : \tint}
  \and
  \inferrule*[right=I-Lam]
    {\Gamma, x{:}\tau_1 \vdash_{\!I} e : \tau_2}
    {\Gamma \vdash_{\!I} \lambda x.\,e : \tau_1 \to \tau_2}
  \and
  \inferrule*[right=I-App]
    {\Gamma \vdash_{\!I} e_1 : \tau_1 \to \tau_2 \\
     \Gamma \vdash_{\!I} e_2 : \tau_1}
    {\Gamma \vdash_{\!I} e_1\,e_2 : \tau_2}
  \and
  \inferrule*[right=I-Let]
    {\Gamma \vdash_{\!I} e_1 : \tau_1 \\
     \gen(\Gamma,\tau_1) \ni \sigma \\
     \Gamma, x{:}\sigma \vdash_{\!I} e_2 : \tau_2}
    {\Gamma \vdash_{\!I} \mathtt{let}\;x = e_1\;\mathtt{in}\;e_2 : \tau_2}
\end{mathpar}

\textsc{I-App} models the situation after unification has succeeded: the
function is already known to have an arrow type whose domain matches the
argument.  \textsc{I-Lam} likewise guesses $\tau_1$ rather than solving for it.
No unification algorithm appears in the development.

\section{Inference soundness: the proof}

\begin{theorem}[\texttt{hm\_infer\_sound}]
  If $\Gamma \vdash_{\!I} e : \tau$ then $\Gamma \vdash e : \tau$.
\end{theorem}

\begin{proof}[Shape of the proof]
Induction on the derivation of $\Gamma \vdash_{\!I} e : \tau$, five cases, one
per constructor.  Each case applies the declarative rule of the same name and
discharges its premises with the ones already at hand.

\begin{center}
\begin{tabular}{lll}
  case & rule applied & premises supplied by \\
  \hline
  \textsc{I-Var} & \textsc{T-Var} & the same two premises, carried over \\
  \textsc{I-Int} & \textsc{T-Int} & none \\
  \textsc{I-Lam} & \textsc{T-Lam} & the induction hypothesis \\
  \textsc{I-App} & \textsc{T-App} & two induction hypotheses \\
  \textsc{I-Let} & \textsc{T-Let} & two induction hypotheses and the
                                    $\gen$ premise \\
\end{tabular}
\end{center}

The Coq proof is correspondingly five tactic lines.

\begin{lstlisting}
Proof.
  intros G e T H. induction H.
  - eapply TyVar; eauto.
  - constructor.
  - apply TyLam. exact IHhm_infer.
  - eapply TyApp; eauto.
  - eapply TyLet; eauto.
Qed.
\end{lstlisting}
\end{proof}

\subsection{Checked examples}

Two programs are derived in the algorithmic system and transported by the
theorem.
\[
  \mathtt{let}\;\mathit{id} = \lambda x.\,x\;\mathtt{in}\;\mathit{id}\;42
  \;:\; \tint
\]
The derivation generalizes the binding to
$\forall[\alpha_0].\,(\tint \to \tint)$ and then uses it at
$\tint \to \tint$ --- the body of the scheme, which by
Remark~\ref{rem:inst} is its only instance.  The second example assumes
$+ : \tint \to \tint \to \tint$ in the environment and derives
\[
  \lambda x.\,(+\;x)\;1 \;:\; \tint \to \tint .
\]

% =====================================================================
\section{Type safety: syntax}
\label{sec:safe-syntax}

\subsection{Grammar}

\[
\begin{array}{lcll}
  t & ::= & x \mid n \mid t_1 + t_2 \mid \lambda x{:}\tau.\,t \mid t_1\,t_2
            \mid \mathtt{let}\;x = t_1\;\mathtt{in}\;t_2
    & \text{terms}\\[2pt]
  v & ::= & n \mid \lambda x{:}\tau.\,t
    & \text{values}\\[2pt]
  \tau & ::= & \tint \mid \tau_1 \to \tau_2
    & \text{types}
\end{array}
\]

\begin{lstlisting}
Inductive ty : Type :=
| TInt : ty
| TArrow : ty -> ty -> ty.

Inductive tm : Type :=
| TVar : nat -> tm       | TIntLit : nat -> tm
| TAdd : tm -> tm -> tm
| TLam : nat -> ty -> tm -> tm
| TApp : tm -> tm -> tm
| TLet : nat -> tm -> tm -> tm.
\end{lstlisting}

Two differences from Section~\ref{sec:sound-syntax} are worth stating
explicitly, because the two files are easily read as one development.

\begin{remark}[The type language has no variables]
\label{rem:notvars}
$\tau ::= \tint \mid \tau_1 \to \tau_2$.  There is no \texttt{TVar}
constructor, no scheme and no quantifier.  Environments map variables to
\emph{types}, not to schemes.
\end{remark}

\begin{remark}[Lambdas are annotated and \texttt{let} is monomorphic]
\label{rem:mono}
$\lambda x{:}\tau.\,t$ carries its argument type, and \textsc{S-Let} below
binds $x$ at the single type given to $t_1$.  There is no generalization step.
\end{remark}

Together these say that the language of this file is the simply typed lambda
calculus with \texttt{let} and addition.

\section{Type safety: rules}

\subsection{Typing}

\begin{mathpar}
  \inferrule*[right=S-Var]
    {\Gamma(x) = \tau}
    {\Gamma \vdash x : \tau}
  \and
  \inferrule*[right=S-Int]
    { }
    {\Gamma \vdash n : \tint}
  \and
  \inferrule*[right=S-Add]
    {\Gamma \vdash t_1 : \tint \\ \Gamma \vdash t_2 : \tint}
    {\Gamma \vdash t_1 + t_2 : \tint}
  \and
  \inferrule*[right=S-Lam]
    {\Gamma, x{:}\tau_1 \vdash t : \tau_2}
    {\Gamma \vdash \lambda x{:}\tau_1.\,t : \tau_1 \to \tau_2}
  \and
  \inferrule*[right=S-App]
    {\Gamma \vdash t_1 : \tau_1 \to \tau_2 \\ \Gamma \vdash t_2 : \tau_1}
    {\Gamma \vdash t_1\,t_2 : \tau_2}
  \and
  \inferrule*[right=S-Let]
    {\Gamma \vdash t_1 : \tau_1 \\ \Gamma, x{:}\tau_1 \vdash t_2 : \tau_2}
    {\Gamma \vdash \mathtt{let}\;x = t_1\;\mathtt{in}\;t_2 : \tau_2}
\end{mathpar}

\subsection{Substitution}

$\subst{t}{x}{s}$ replaces free occurrences of $x$ by $s$, stopping at a
binder for $x$.

\[
\begin{array}{lcl}
  \subst{y}{x}{s} &=& \text{$s$ if $y = x$, else $y$} \\
  \subst{n}{x}{s} &=& n \\
  \subst{(t_1 + t_2)}{x}{s} &=& \subst{t_1}{x}{s} + \subst{t_2}{x}{s} \\
  \subst{(\lambda y{:}\tau.\,t)}{x}{s} &=&
     \text{$\lambda y{:}\tau.\,t$ if $y = x$, else
           $\lambda y{:}\tau.\,\subst{t}{x}{s}$} \\
  \subst{(t_1\,t_2)}{x}{s} &=& \subst{t_1}{x}{s}\;\subst{t_2}{x}{s} \\
  \subst{(\mathtt{let}\;y = t_1\;\mathtt{in}\;t_2)}{x}{s} &=&
     \mathtt{let}\;y = \subst{t_1}{x}{s}\;\mathtt{in}\;
     (\text{$t_2$ if $y = x$, else $\subst{t_2}{x}{s}$})
\end{array}
\]

\subsection{Operational semantics}

Call by value, left to right.

\begin{mathpar}
  \inferrule*[right=E-Add]
    { }
    {n + m \step n{+}m}
  \and
  \inferrule*[right=E-Add1]
    {t_1 \step t_1'}
    {t_1 + t_2 \step t_1' + t_2}
  \and
  \inferrule*[right=E-Add2]
    {t \step t'}
    {v + t \step v + t'}
  \and
  \inferrule*[right=E-Beta]
    { }
    {(\lambda x{:}\tau.\,t)\,v \step \subst{t}{x}{v}}
  \and
  \inferrule*[right=E-App1]
    {t_1 \step t_1'}
    {t_1\,t_2 \step t_1'\,t_2}
  \and
  \inferrule*[right=E-App2]
    {t \step t'}
    {v\,t \step v\,t'}
  \and
  \inferrule*[right=E-Let]
    { }
    {\mathtt{let}\;x = v\;\mathtt{in}\;t \step \subst{t}{x}{v}}
  \and
  \inferrule*[right=E-Let1]
    {t_1 \step t_1'}
    {\mathtt{let}\;x = t_1\;\mathtt{in}\;t_2 \step
     \mathtt{let}\;x = t_1'\;\mathtt{in}\;t_2}
\end{mathpar}

\section{Type safety: the proof}

\subsection{Environment lemmas}

Four facts about \texttt{extend}, each by case analysis on $\mathtt{Nat.eqb}$:

\[
\begin{array}{ll}
  \texttt{extend\_eq} & (\Gamma, x{:}\tau)(x) = \tau \\
  \texttt{extend\_neq} & x \neq y \Rightarrow (\Gamma, x{:}\tau)(y) = \Gamma(y) \\
  \texttt{extend\_shadow} & \Gamma, x{:}\tau_1, x{:}\tau_2 = \Gamma, x{:}\tau_2 \\
  \texttt{extend\_permute} & x \neq y \Rightarrow
     \Gamma, x{:}\tau_x, y{:}\tau_y = \Gamma, y{:}\tau_y, x{:}\tau_x
\end{array}
\]

On these sit \texttt{context\_invariance} (pointwise-equal environments type
the same terms) and \texttt{weakening} (typing survives extending the
environment), both by induction on the typing derivation, with
\texttt{weaken\_empty} as the special case from $\emptyset$.

\subsection{Substitution}

\begin{lemma}[\texttt{substitution\_preserves\_typing}]
  If $\Gamma, x{:}\upsilon \vdash t : \tau$ and $\vdash v : \upsilon$ then
  $\Gamma \vdash \subst{t}{x}{v} : \tau$.
\end{lemma}

\begin{proof}[Shape of the proof]
Induction on the term $t$, inverting the typing derivation in each case.  Six
cases; three carry the argument.

\emph{Variable.}  Split on whether the variable is $x$.  If it is,
\texttt{extend\_eq} identifies its type as $\upsilon$ and
\texttt{weaken\_empty} lifts $\vdash v : \upsilon$ to $\Gamma$.  If it is not,
\texttt{extend\_neq} discards the binding and \textsc{S-Var} rebuilds the
derivation.

\emph{Lambda.}  For $\lambda y{:}\tau_1.\,t'$ the substitution stops when
$y = x$; that case is closed by \texttt{context\_invariance} together with
\texttt{extend\_shadow}, since the inner binding of $x$ hides the outer one.
When $y \neq x$ the two bindings are exchanged by \texttt{extend\_permute} and
the induction hypothesis applies.

\emph{Let.}  The same shadow/permute split for the body; the bound expression
is handled by the induction hypothesis directly.
\end{proof}

\subsection{Canonical forms}

By inversion on $\mathrm{value}(v)$ and then on the typing derivation:
\[
  \mathrm{value}(v),\; \vdash v : \tint
    \;\Longrightarrow\; v = n
  \qquad
  \mathrm{value}(v),\; \vdash v : \tau_1 \to \tau_2
    \;\Longrightarrow\; v = \lambda x{:}\tau_1.\,t
\]
The remaining combinations are refuted because \textsc{S-Int} and
\textsc{S-Lam} build different type constructors.

\subsection{Progress}

\begin{theorem}[\texttt{progress}]
  If $\vdash t : \tau$ then $t$ is a value or $t \step t'$ for some $t'$.
\end{theorem}

\begin{proof}[Shape of the proof]
Induction on the typing derivation with the environment remembered as
$\emptyset$.

\textsc{S-Var} is vacuous: $\emptyset$ yields no binding, so its premise is
contradictory.  \textsc{S-Int} and \textsc{S-Lam} produce values.

\textsc{S-Add}: if both operands are values, canonical forms make them integer
literals and \textsc{E-Add} fires; otherwise whichever operand steps is
propagated by \textsc{E-Add1} or \textsc{E-Add2}.  \textsc{S-App} runs the same
way, with the canonical form at arrow type turning the function into a lambda
so that \textsc{E-Beta} applies.  \textsc{S-Let}: if $t_1$ is a value then
\textsc{E-Let} fires, else \textsc{E-Let1} propagates.
\end{proof}

\subsection{Preservation}

\begin{theorem}[\texttt{preservation}]
  If $\vdash t : \tau$ and $t \step t'$ then $\vdash t' : \tau$.
\end{theorem}

\begin{proof}[Shape of the proof]
Induction on the step derivation, inverting the typing derivation in each case.
The congruence rules rebuild the same typing rule around the induction
hypothesis, and \textsc{E-Add} is closed by \textsc{S-Int}.  The two
substituting rules are where the work is.  \textsc{E-Beta} inverts
\textsc{S-App} and then \textsc{S-Lam} to expose
$\Gamma, x{:}\tau_1 \vdash t : \tau_2$ and $\vdash v : \tau_1$, which are the
two hypotheses of the substitution lemma.  \textsc{E-Let} inverts
\textsc{S-Let} and applies the same lemma.
\end{proof}

\texttt{type\_safety} is stated as a corollary and is definitionally
\texttt{preservation}.

\subsection{Checked example}

\[
  \mathtt{let}\;f = \lambda x{:}\tint.\,x + 1\;\mathtt{in}\;f\;41
\]
is proved to have type $\tint$, stepped once by \textsc{E-Let}, and re-typed
through \texttt{preservation}.

% =====================================================================
\section{How to read these results}
\label{sec:reading}

Both files compile and every stated result is closed with \texttt{Qed}.  What
they assert, however, is narrower than their names suggest, and this section
says how much.

\subsection{The soundness theorem is a restatement}

Compare the two systems rule by rule.  \textsc{I-Var} and \textsc{T-Var} have
the same two premises; \textsc{I-Int} and \textsc{T-Int} have none;
\textsc{I-Lam}, \textsc{I-App} and \textsc{I-Let} agree with \textsc{T-Lam},
\textsc{T-App} and \textsc{T-Let} premise for premise.  The two relations are
therefore \emph{the same relation}, presented twice, and
\texttt{hm\_infer\_sound} is proved by sending each constructor to its twin.

This is not a defect in the proof --- the proof is correct --- but it means the
theorem carries no information about type inference.  In particular:

\begin{itemize}
  \item no unification algorithm is defined, so nothing is proved about one;
  \item instantiation is the identity (Remark~\ref{rem:inst}), so the theorem
  says nothing about using a scheme at several types;
  \item generalization is unconstrained (Remark~\ref{rem:gen}), so the
  Damas--Milner side condition is neither assumed nor established.
\end{itemize}

To make the statement carry content one would let \textsc{Inst} substitute for
the quantified variables, give \textsc{Gen} the condition
$\bar\alpha \cap \ftv(\Gamma) = \emptyset$ --- which requires environments to
be association lists, so that $\ftv(\Gamma)$ can be computed --- and let the
algorithmic system compute a particular scheme rather than choose one.  The two
systems would then differ, and soundness would have to prove that what the
algorithm computes satisfies what the declarative system demands.

\subsection{The safety theorem is about a simply typed language}

By Remarks~\ref{rem:notvars} and~\ref{rem:mono}, \texttt{HMTypeSafety.v} has no
type variables, no schemes and no polymorphic \texttt{let}.  Progress and
preservation are proved, and proved properly --- the substitution lemma,
context invariance, weakening and canonical forms are all there --- but for the
simply typed lambda calculus with \texttt{let} and addition.  Nothing
Hindley--Milner occurs in it.

The two files also do not connect: the type language of
Section~\ref{sec:safe-syntax} is not that of Section~\ref{sec:sound-syntax}, so
a program typed by \texttt{hm\_infer\_sound} is not thereby a program to which
\texttt{progress} applies.

\subsection{Not covered}

\begin{itemize}
  \item a concrete unification algorithm and its correctness;
  \item the occurs check;
  \item a principal type theorem;
  \item polymorphism in the language that is proved safe;
  \item recursion: neither file has \texttt{let rec}, so nothing here bears on
  recursive definitions;
  \item lists, pairs, pattern matching, references, exceptions;
  \item any correspondence with the OCaml prover in this repository
  (\texttt{w.ml}, \texttt{typing\_thm.ml}, \texttt{prove.ml} and the rest).
  The Coq development and the OCaml code were written separately, and no
  theorem relates them.
\end{itemize}

\section{Verification}

\begin{lstlisting}
coqc HMSoundness.v
coqc HMTypeSafety.v
\end{lstlisting}

Both compile.  Neither file contains \texttt{Admitted}, \texttt{admit} or
\texttt{Axiom}, so every stated result is fully proved from the definitions
above --- which is precisely why Section~\ref{sec:reading} matters.  The proofs
are sound about the systems they define, and those systems are smaller than
their names.

The report is built with:

\begin{lstlisting}
pdflatex -interaction=nonstopmode hm_safety_report.tex   # run twice
\end{lstlisting}

\end{document}
